\section{Generator and evaluation protocol}
\label{app:protocol}

\paragraph{Data.}
BCI Competition IV-2a \citep{brunner2008bci,tangermann2012review}: nine subjects, four-class motor imagery, 22 EEG channels at 250~Hz.
The recordings are cut into 750-sample windows, giving 2592 trials, which are the units of the hierarchical generator.

\paragraph{Generator training.}
A shallow piecewise-linear recurrent neural network \citep{brenner2024shPLRNN} with $M$ latent states, hidden width 20 and a linear observation model $x = Wz$, $W \in \mathbb{R}^{22 \times M}$, is trained with generalised teacher forcing \citep{hess2023gtf} (learning rate $10^{-4}$ for the shared parameters and $10^{-3}$ for the trial-specific ones, batch size 64, sequence length 50, 500 epochs).
One generator is trained per $M \in \{2, 4, 8, 16, 22, 32\}$ with the same recipe, and each supplies its own $W$.

\paragraph{Source trajectories.}
The deterministic free run of a trained generator contracts to a fixed point, so the sources are its stochastic free run $z_t = f(z_{t-1}) + \epsilon_t$, $\epsilon_t \sim \mathcal{N}(0, \mathrm{diag}(\sigma^2))$, where $\sigma$ is the per-latent median of the generator's one-step prediction residual over all 2592 training trials.
Each seed draws one trial's parameters and starts from that trial's first real frame; 500 burn-in steps are discarded, $T = 2000$ steps are kept, each latent is standardised to unit variance (the scores are invariant to this), the sources are projected through $W$ and the configured nuisance is added.

\paragraph{Held-out protocol.}
For every cell two trajectories are drawn from the same generator with the same trial parameters and initial frame (generator seeds $s$ and $s + 50000$).
One draw of the condition randomness (segment-wise modulation gains, coupling matrix, $W$, noise type and noise power, the latter set from the training trajectory's clean signal power) is applied to both.
The training trajectory is standardised with its own statistics and the test trajectory with the training statistics.
Every method is fit on the training trajectory and applied to the test one; permutation matching and the canonical-correlation map are fit on the training trajectory and applied unchanged to the test one.

\paragraph{Segments and modulation.}
The $T = 2000$ steps are cut into 20 equal segments, which are the labels given to TCL, iVAE and segment-conditioned CEBRA.
Unless the modulation axis says otherwise, each source is rescaled per segment by an independent $\mathrm{Exp}(1)$ variance (source-specific profiles).
Along axis (d) the per-source log-variance profile is $((1-a)\,g_0 + a\,g_i)/\sqrt{(1-a)^2 + a^2}$ with $g_0$ a profile common to all sources and $g_i$ a source-specific one, for $a \in \{1, 0.75, 0.5, 0.35, 0.2, 0.1, 0\}$; $a = 1$ is the source-specific law and $a = 0$ a single shared gain (rank-one modulation matrix).
Cells with $r \in \{1, 2, 3, 4\}$ distinct profiles shared among the four sources fall on the same curve (Table~\ref{tab:mod}).

\paragraph{Other axes.}
Pink noise is generated by filtering Gaussian white noise with a $1/f$ amplitude spectrum and scaled to the configured signal-to-noise ratio ($\{30, 15, 10, 5, 0\}$~dB); the remaining axes use 20~dB Gaussian sensor noise.
Mixing nonlinearity applies $\phi \in \{\mathrm{id},\, x \mapsto x^3,\, \tanh\}$ pointwise after $W$.
Source coupling applies $z_t \leftarrow (1-\gamma)\, z_t + \gamma\, A z_{t-1}$ with $A$ a fixed random matrix of unit spectral norm, $\gamma \in \{0, 0.1, 0.3, 0.6, 0.9\}$, after the segment rescaling and before mixing.
At $M = 32 > D$ the linear methods return 22 components and are scored over the 22 matched pairs.

\paragraph{Seeds.}
Twenty seeds for PCA, FastICA and TCL; five for iVAE, PCL, demeaned TCL, axis (d) and $M = 32$; three for CEBRA.
Each seed draws a different trial's generator parameters and a different fitting seed.

\section{Methods and implementation}
\label{app:methods}

\paragraph{PCA, FastICA.} \texttt{scikit-learn} with default tolerances, $M$ components.
\paragraph{TCL.} Per-timestep MLP feature extractor (three layers, hidden width 64, LeakyReLU) and a multinomial logistic regression over the 20 segments; Adam, learning rate $10^{-3}$, 200 epochs; the $M$-dimensional features are the components \citep{hyvarinen2016tcl}.
\textit{Demeaned TCL} removes the per-segment, per-channel mean from the training and the test observations before fitting; this removes only power below the segment rate and keeps 91\% of the source power.
\paragraph{iVAE.} Encoder and decoder MLPs (three layers, hidden width 64), a conditional Gaussian prior factorised over latent dimensions with the segment index as auxiliary variable, per-channel standardised inputs and the evidence lower bound under a unit-variance Gaussian likelihood; Adam, 200 epochs; the posterior mean is the component \citep{khemakhem2020ivae}.
\paragraph{PCL.} The TCL encoder with the per-component dependence model of \citet{hyvarinen2017pcl}, trained for 200 epochs to discriminate consecutive pairs $(x_t, x_{t-1})$ from time-permuted pairs; scored like the source-estimating methods.
PCL has no positive control on the generator's sources, whose innovations are Gaussian; on AR(1) sources with Laplacian innovations the same implementation scores 0.88 (stationary) and 0.90 (modulated), on Gaussian AR(1) sources 0.71, the rotation level.
\paragraph{CEBRA.} CEBRA 0.6.1 \citep{schneider2023cebra}, per-timestep MLP with a Euclidean (non-normalised) output of dimension $M$, 500 iterations, time-contrastive positives at offsets up to 10 samples (CEBRA-Time).
In every cell its features are energy-like (canonical correlation with $|s|$ well above that with $s$), so they are scored as TCL's.
A segment-conditioned variant gives the same profile.
\paragraph{Hardware.} All fits run on CPU (Apple M4); generator training on Apple MPS.

\section{Scores, chance level and rotation reference}
\label{app:scores}

\paragraph{Permutation MCC.} Mean absolute correlation between components and signed sources after solving the linear assignment problem on the absolute correlation matrix \citep{hyvarinen2016tcl}; the assignment is fit on the training trajectory.
\paragraph{Linear-invariant MCC.} Under segment-wise variance modulation of Gaussian sources the sufficient statistic is $q(s) = s^2$, and TCL identifies $q(s)$ only up to a linear map (Appendix~\ref{app:tcl}).
TCL and CEBRA are therefore scored by the mean canonical correlation between their features and $|s|$, with the canonical map fit on the training trajectory.
$|s|$ is a lighter-tailed monotone surrogate of $s^2$; scoring against $s^2$ instead lowers the values by $0.01$--$0.08$ without changing which conditions recover and which fail (Table~\ref{tab:mod}, last column).
\paragraph{Chance level.} Both scores are maximised over a matching, so an estimator carrying no information does not score zero.
For each cell the true sources are drawn as in the sweep and scored against surrogate components that are statistically matched to them but independent of them: an independent draw from the same generator, and a circular time shift of the sources themselves, which preserves every autocorrelation and destroys only the alignment.
The band in the figures is the 95th percentile over 20 surrogate repeats, maximised over both constructions and both scores.
\paragraph{Rotation reference.} The permutation MCC of the true sources under a Haar-random orthogonal rotation (20 repeats, same matching): $0.92$, $0.77$, $0.66$, $0.54$ and $0.47$ at $M = 2, 4, 8, 16, 22$.
A method that recovers the source subspace without separating it scores at this level; it is specific to each $M$ and to the sources' correlation structure, and it does not apply to the linear-invariant score, which is unchanged by any rotation.

\section{Full results}
\label{app:tables}

\begin{table}[h]\centering\small
\caption{Held-out recovery on the benchmark of Figure~\ref{fig:axes}. PCA, FastICA and TCL: mean over 20 seeds with bootstrap 95\% interval; iVAE, PCL and demeaned TCL: mean over 5 seeds. Permutation MCC against the signed sources for PCA, FastICA, iVAE and PCL (rotation reference in the last column); CCA-MCC against $|s|$ for TCL. Null: 95th percentile of the chance level over surrogate constructions and scores. The modulation block was run at 15~dB $1/f$, so its source-specific cell is the 15~dB cell of the first block; the modulation axis of Figure~\ref{fig:axes}d at the clean operating point is in Table~\ref{tab:mod}.}
\label{tab:full}
\begin{tabular}{llcccccccc}\toprule
axis & cell & PCA & FastICA & iVAE & PCL & TCL & TCL demeaned & null & rotation \\ \midrule
$1/f$ SNR (dB) & 30 & 0.71 [0.68, 0.74] & 0.98 [0.96, 0.99] & 0.92 & 0.58 & 0.81 [0.78, 0.84] & 0.89 & 0.05 & 0.77 \\
 & 15 & 0.70 [0.68, 0.73] & 0.97 [0.96, 0.98] & 0.92 & 0.31 & 0.39 [0.35, 0.43] & 0.86 & 0.05 & 0.77 \\
 & 10 & 0.70 [0.67, 0.73] & 0.96 [0.95, 0.98] & 0.89 & 0.28 & 0.27 [0.24, 0.29] & 0.78 & 0.05 & 0.77 \\
 & 5 & 0.68 [0.66, 0.71] & 0.94 [0.92, 0.95] & 0.84 & 0.19 & 0.19 [0.17, 0.21] & 0.58 & 0.05 & 0.77 \\
 & 0 & 0.63 [0.59, 0.66] & 0.86 [0.83, 0.88] & 0.43 & 0.16 & 0.13 [0.11, 0.14] & 0.39 & 0.05 & 0.77 \\
\midrule
mixing & linear & 0.71 [0.68, 0.74] & 0.97 [0.96, 0.99] & 0.93 & 0.55 & 0.89 [0.85, 0.92] & -- & 0.05 & 0.77 \\
 & cube & 0.67 [0.62, 0.72] & 0.80 [0.79, 0.82] & 0.71 & 0.44 & 0.72 [0.70, 0.75] & -- & 0.05 & 0.77 \\
 & tanh & 0.61 [0.59, 0.63] & 0.64 [0.63, 0.66] & 0.87 & 0.61 & 0.77 [0.74, 0.80] & -- & 0.05 & 0.77 \\
\midrule
coupling $\gamma$ & 0.0 & 0.71 [0.68, 0.74] & 0.97 [0.96, 0.99] & 0.93 & 0.55 & 0.89 [0.85, 0.92] & -- & 0.05 & 0.77 \\
 & 0.1 & 0.71 [0.68, 0.74] & 0.97 [0.96, 0.99] & 0.94 & 0.54 & 0.89 [0.85, 0.92] & -- & 0.05 & 0.77 \\
 & 0.3 & 0.72 [0.68, 0.75] & 0.97 [0.96, 0.98] & 0.92 & 0.65 & 0.87 [0.84, 0.91] & -- & 0.06 & 0.77 \\
 & 0.6 & 0.74 [0.71, 0.78] & 0.97 [0.96, 0.98] & 0.87 & 0.61 & 0.83 [0.79, 0.87] & -- & 0.06 & 0.77 \\
 & 0.9 & 0.73 [0.69, 0.77] & 0.93 [0.91, 0.95] & 0.81 & 0.55 & 0.75 [0.71, 0.79] & -- & 0.09 & 0.77 \\
\midrule
modulation (15~dB $1/f$) & source-specific & 0.70 [0.68, 0.73] & 0.97 [0.96, 0.98] & 0.92 & 0.31 & 0.39 [0.35, 0.43] & 0.86 & 0.05 & 0.77 \\
 & shared gain & 0.66 [0.64, 0.68] & 0.79 [0.76, 0.82] & 0.61 & 0.29 & 0.28 [0.27, 0.29] & 0.33 & 0.04 & 0.78 \\
 & none (stationary) & 0.64 [0.61, 0.66] & 0.72 [0.68, 0.76] & 0.59 & 0.36 & 0.05 [0.05, 0.06] & 0.17 & 0.06 & 0.78 \\
\midrule
$M$ ($D = 22$) & 2 & 0.90 [0.85, 0.93] & 0.98 [0.98, 0.99] & 0.98 & 0.84 & 0.90 [0.85, 0.94] & 0.89 & 0.08 & 0.92 \\
 & 4 & 0.71 [0.68, 0.74] & 0.97 [0.96, 0.99] & 0.93 & 0.55 & 0.89 [0.85, 0.92] & 0.89 & 0.05 & 0.77 \\
 & 8 & 0.61 [0.60, 0.63] & 0.96 [0.95, 0.97] & 0.64 & 0.51 & 0.81 [0.79, 0.82] & 0.77 & 0.05 & 0.66 \\
 & 16 & 0.50 [0.49, 0.51] & 0.95 [0.94, 0.96] & 0.32 & 0.40 & 0.57 [0.56, 0.58] & 0.55 & 0.03 & 0.54 \\
 & 22 & 0.44 [0.43, 0.45] & 0.90 [0.88, 0.92] & 0.28 & 0.37 & 0.41 [0.41, 0.42] & 0.41 & 0.03 & 0.47 \\
\bottomrule
\end{tabular}
\end{table}

\begin{table}[h]\centering\small
\caption{Modulation-diversity axis (Figure~\ref{fig:axes}d and Figure~\ref{fig:real}a; $M = 4$, 20~dB Gaussian noise, 5 seeds): held-out recovery against the observed co-modulation fraction of the segment-wise modulation. $a$ is the source-specific weight of the profile, $r$ the number of distinct profiles. The nine BCI IV-2a subjects have shares of 0.22--0.76 at 1--40~Hz and 0.24--0.42 at 8--30~Hz.}
\label{tab:mod}
\begin{tabular}{lcccccccc}\toprule
cell & co-modulation fraction & rank & PCA & FastICA & iVAE & PCL & TCL & TCL vs $s^2$ \\ \midrule
$a = 1$ & 0.42 & 4 & 0.71 & 0.97 & 0.92 & 0.57 & 0.89 & 0.81 \\
$a = 0.75$ & 0.39 & 4 & 0.70 & 0.96 & 0.93 & 0.51 & 0.85 & 0.78 \\
$a = 0.5$ & 0.61 & 4 & 0.68 & 0.95 & 0.86 & 0.52 & 0.84 & 0.78 \\
$a = 0.35$ & 0.82 & 4 & 0.66 & 0.94 & 0.76 & 0.49 & 0.72 & 0.67 \\
$a = 0.2$ & 0.95 & 4 & 0.63 & 0.87 & 0.62 & 0.53 & 0.52 & 0.50 \\
$a = 0.1$ & 0.99 & 4 & 0.63 & 0.75 & 0.59 & 0.49 & 0.40 & 0.38 \\
$a = 0$ & 0.99 & 1 & 0.64 & 0.75 & 0.59 & 0.46 & 0.32 & 0.31 \\
$r = 3$ & 0.59 & 3 & 0.69 & 0.89 & 0.81 & 0.47 & 0.74 & 0.68 \\
$r = 2$ & 0.63 & 2 & 0.66 & 0.89 & 0.72 & 0.45 & 0.53 & 0.48 \\
$r = 1$ & 0.99 & 1 & 0.64 & 0.74 & 0.58 & 0.45 & 0.35 & 0.32 \\
\bottomrule\end{tabular}\end{table}

\section{Robustness of the benchmark}
\label{app:robust}

\begin{figure}[h]
  \centering
  \includegraphics[width=0.78\textwidth]{figR_robustness.pdf}
  \caption{Robustness of Figure~\ref{fig:axes}. (a) Held-out recovery against the $1/f$ signal-to-noise ratio without preprocessing (solid) and after per-segment demeaning of the observations (dashed). (b) TCL against the length $T$ of the training trajectory at the stressed cells, 20 segments (lines) and 320 segments at $T = 32000$ (hollow). (c) Recovery with an additional nuisance component that is not among the sources, at $-10$ to $+10$~dB relative to the signal, with its variance modulated across segments like the sources' (solid) or constant (dashed). (d) Recovery with a sphere-model leadfield in place of the learned mixing, one point per cell of Figure~\ref{fig:axes}. Means over 5 seeds, bootstrap 95\% intervals.}
  \label{fig:robust}
\end{figure}

\paragraph{Preprocessing.}
Removing the per-segment mean of the observations restores TCL on the $1/f$ axis to its clean level down to 10~dB and leaves the other methods within a few hundredths of their values (Figure~\ref{fig:robust}a); its held-out segment accuracy rises at the same time, showing that the classifier had been fitting drift.
Zero-phase high-pass filters at one to four times the segment rate have the same effect on TCL and remove 18--35\% of the source power, which costs the other methods.
The shared-gain and stationary cells stay failed after any preprocessing.

\paragraph{Data quantity.}
With the same generator and mixing, TCL's score at $M = 22$ rises monotonically with the length of the training trajectory and with the number of segments, up to the clean level at $T = 32000$ with 320 segments (Figure~\ref{fig:robust}b); the linear methods are already saturated.
The shared-gain and stationary cells do not move with data, as the rank argument predicts, and the $1/f$ cells become worse with more segments, which give the classifier more drift signatures to memorise.

\paragraph{Group energy under shared profiles.}
With $r$ distinct modulation profiles shared among the four sources, the energies $\sum_{i \in g} s_i^2$ of the groups sharing a profile are decoded from the frozen TCL features with held-out $R^2$ above $0.8$ at every $r$, while the squares of individual sources within a group are decoded only up to the ceiling implied by the rotation ambiguity; rank deficiency coarsens the identified statistics to the invariants of the ambiguity group rather than destroying them.

\paragraph{Nuisance component.}
One additional white Gaussian signal on a random unit-norm spatial pattern, not among the sources, costs TCL about $0.13$ as soon as its variance is modulated across segments, already at a tenth of the signal power and no more at ten times it: one of TCL's features tracks the nuisance envelope (Figure~\ref{fig:robust}c).
A stationary nuisance of the same power costs TCL $0.02$--$0.06$.
The other methods lose one component to the nuisance once its power exceeds $-5$~dB, independent of modulation.

\paragraph{Leadfield.}
Replacing the learned $W$ by a physics leadfield (MNE three-layer sphere model fitted to the 22 BCI IV-2a electrodes \citep{gramfort2014mne}, $M$ current dipoles at random positions and orientations, rescaled to the learned $W$'s mean row norm) reproduces the ordering of every axis of Figure~\ref{fig:axes} (Figure~\ref{fig:robust}d).
The leadfield's condition number is $10^3$ at $M = 16$ and $10^4$ at $M = 22$, so the decline with $M$ is steeper for every method than with the learned mixing.

\paragraph{Capacity.}
A TCL network with four times the hidden width or five times the epochs changes the stressed cells by at most $0.04$ and lowers the clean cell, so the reported setting is the strongest of those tested.

\section{Generator realism}
\label{app:realism}

\begin{figure}[h]
  \centering
  \includegraphics[width=0.6\textwidth]{figA_generator_realism.pdf}
  \caption{(a) Normalised power spectra of real BCI IV-2a EEG (per subject) and of the generator's synthetic EEG $Wz$ (per seed) for the recipe of Appendix~\ref{app:protocol} and for a recipe with a generalised-teacher-forcing floor of $0.2$ and hidden width 64; (b) the corresponding latent spectra. Thin lines: individual subjects or seeds; thick: median.}
  \label{fig:realism}
\end{figure}

The latents that the generator encodes from real EEG have a $1/f$ spectrum with an alpha peak; its stochastic free run under the recipe of Appendix~\ref{app:protocol} is flatter (Figure~\ref{fig:realism}).
Training with a generalised-teacher-forcing floor of $0.2$ and hidden width 64 gives free runs whose spectral slope, alpha excess and lag-one autocorrelation match the encoded latents.
Rerunning the full benchmark on generators trained with this recipe reproduces Figure~\ref{fig:axes} qualitatively (Figure~\ref{fig:axes_realistic}): TCL fails on the stationary and shared-gain cells and recovers along the $1/f$ axis only after demeaning, PCL survives the stationary cell, and every method falls with $M$.
These sources are more strongly inter-correlated than the main benchmark's, so the rotation reference and the chance band are higher at $M = 4$ (0.83 and 0.13), and the free runs are EEG-like for $M \le 8$.

\begin{figure}[h]
  \centering
  \includegraphics[width=\textwidth]{fig1_axes_realistic.pdf}
  \caption{Figure~\ref{fig:axes} on the generators of Appendix~\ref{app:realism} (5 seeds; panel (d) uses the three categorical modulation cells).}
  \label{fig:axes_realistic}
\end{figure}

\section{Real-EEG protocols}
\label{app:real}

\subsection{Observable properties of the segment-wise modulation}
\label{app:placement}
\paragraph{Co-modulation fraction.}
The recording is cut into the 20 contiguous segments TCL uses.
Each segment's channel covariance is whitened by the mean covariance and mapped to the tangent space by the matrix logarithm; the 20 tangent vectors are centred and their singular values taken.
The co-modulation fraction is the fraction of the segment-to-segment change that follows the dominant pattern, after subtracting the mean spectrum of 20 stationary surrogates (multivariate phase randomisations of the whole recording, which keep every auto- and cross-spectrum and remove the nonstationarity).
A single shared profile gives a share of one.
The statistic is computed identically on every synthetic cell of axis (d) and on the nine BCI IV-2a subjects (all six motor-imagery runs, 22 channels, 1--40~Hz and 8--30~Hz).
Real subjects have shares of 0.22--0.76 (median 0.39) at 1--40~Hz and 0.24--0.42 (median 0.32) at 8--30~Hz; on the benchmark TCL falls below $0.5$ only at shares of $0.95$ and above.

\paragraph{Magnitude.}
The magnitude of the modulation is the standard deviation, across segments and components, of the per-segment log-variance in dB.
On the benchmark it is the nominal standard deviation of the log-variance profile: $4.3$~dB for the $\mathrm{Exp}(1)$ law of the main sweep, and $2.2$, $1.1$ and $0.4$~dB for the cells of Figure~\ref{fig:real}a, obtained by scaling the log-variance profiles by $0.5$, $0.25$ and $0.1$ with everything else unchanged (Table~\ref{tab:mag}).
On real EEG it is computed on four FastICA components of the 1--40~Hz recording of each subject, cut into 20, 200 or 2000 contiguous segments, with the estimation floor (the same statistic on white noise of the same length) removed in quadrature; the two sides measure the same quantity with different estimators.
Across the nine subjects the median magnitude is $0.4$, $1.0$ and $2.0$~dB at 20, 200 and 2000 segments.

\begin{table}[h]\centering\small
\caption{Held-out recovery against the magnitude of the segment-wise modulation ($M = 4$, 20~dB Gaussian noise, 5 seeds; Figure~\ref{fig:real}a). Source-specific profiles ($a = 1$) for all methods; TCL also at $a = 0.5$ and under one shared gain ($a = 0$).}
\label{tab:mag}
\begin{tabular}{lcccccc}\toprule
magnitude (dB) & PCA & FastICA & iVAE & TCL & TCL, $a = 0.5$ & TCL, shared gain \\ \midrule
0.4 & 0.64 & 0.72 & 0.60 & 0.31 & 0.29 & 0.26 \\
1.1 & 0.65 & 0.81 & 0.65 & 0.47 & 0.45 & 0.32 \\
2.2 & 0.68 & 0.95 & 0.75 & 0.75 & 0.65 & 0.33 \\
4.3 & 0.71 & 0.97 & 0.92 & 0.89 & 0.84 & 0.32 \\
\bottomrule\end{tabular}\end{table}

\subsection{Scoring unsupervised components by the cue}
\label{app:cue}
Per subject, the six motor-imagery runs are band-passed to 8--30~Hz and resampled to 125~Hz.
Each unsupervised method ($M = 4$ or $8$ components) is fit on the continuous signal of runs 1--3 without labels and applied to runs 4--6.
Per trial, the log-variance of each component in the window 0.5--4~s after the cue is the feature; a four-class shrinkage LDA with five-fold cross-validation gives the decoding accuracy, and 20 label permutations its null distribution.
Supervised one-versus-rest common spatial patterns on all 22 channels in the same window give the label-aware ceiling.
Run-to-run agreement is the mean canonical correlation between the components of two fits of the same method on runs 1--3 and 4--6, evaluated on the held-out runs.
Comparisons between methods are paired Wilcoxon signed-rank tests over the nine subjects: at 2000 segments TCL's decoding accuracy differs from PCA, FastICA and PCL at $p \ge 0.055$ for $M = 4$ and $p \ge 0.13$ for $M = 8$, CSP exceeds TCL at $p = 0.004$, and TCL's run-to-run agreement is below every other method at every segment length at $p = 0.004$.
TCL is trained with 20, 200 and 2000 contiguous segments of the continuous signal (about 2~min, 12~s and 1.2~s per segment); the other methods do not depend on the segmentation (Table~\ref{tab:cue}).

\begin{table}[h]\centering\small
\caption{Real BCI IV-2a, median over nine subjects (Figure~\ref{fig:real}b,c). Cue decoding accuracy from all components and run-to-run agreement, for $M = 4$ / $M = 8$ components. The label-permutation 95th percentile is 0.31; supervised CSP decodes at 0.63.}
\label{tab:cue}
\begin{tabular}{llcccc}\toprule
 & & \multicolumn{3}{c}{TCL, segments} & \\ \cmidrule(lr){3-5}
 & & 20 & 200 & 2000 & PCA / FastICA / PCL \\ \midrule
cue decoding & $M = 4$ & 0.35 & 0.37 & 0.41 & 0.43 / 0.37 / 0.44 \\
 & $M = 8$ & 0.33 & 0.44 & 0.48 & 0.40 / 0.46 / 0.49 \\
run-to-run agreement & $M = 4$ & 0.51 & 0.60 & 0.61 & 1.00 / 0.97 / 0.92 \\
 & $M = 8$ & 0.48 & 0.57 & 0.70 & 0.99 / 0.95 / 0.85 \\
\bottomrule\end{tabular}\end{table}

\section{TCL identifiability and the rank condition}
\label{app:tcl}

We follow \citet{hyvarinen2016tcl}.
The data are cut into $K$ segments, and a feature extractor $h(x_t; \theta)$ is trained with per-segment logistic-regression weights $w_\tau$ and biases $b_\tau$ to predict the segment index $C_t = \tau$ from $x_t$, $p(C_t = \tau \mid x_t) \propto \exp(w_\tau^\top h(x_t; \theta) + b_\tau)$.
By Bayes' rule the posterior also equals $p_\tau(x_t)\, p(C_t = \tau)$ up to normalisation, where $p_\tau$ is the segment-wise density, so at the optimum
\begin{equation}
  w_\tau^\top h(x_t; \theta) + b_\tau = \log p_\tau(x_t) - \log p_1(x_t) + \mathrm{const}.
  \label{eq:tcl_post}
\end{equation}
Suppose $x = f(s)$ for an invertible $f$ and the sources have a conditionally exponential-family form $\log p_\tau(s_i) = q_i(s_i)\,\lambda_i(\tau) + \mathrm{const}$, independent across components.
Theorem~1 of \citet{hyvarinen2016tcl} further requires (their assumption A3) that the modulation matrix $L$ with entries $L_{\tau i} = \lambda_i(\tau) - \lambda_i(1)$ has full column rank $M$: the sources must change across segments in linearly independent ways.
Substituting into Eq.~\eqref{eq:tcl_post} and matching segment-dependent terms yields $h(x) \approx A\, q(s) + d$ for some matrix $A$ and shift $d$, with $q(s) = (q_1(s_1), \dots, q_M(s_M))$ and $q_i(s_i) = s_i^2$ for Gaussian sources.
If $L$ has rank $r < M$, the right-hand side of Eq.~\eqref{eq:tcl_post} depends on $s$ only through $r$ linear combinations of $q(s)$, and only those are identified.
For a gain shared by all sources, $\lambda_i(\tau) - \lambda_i(1) = c(\tau)\, a_i$, $L$ has rank one and the single identified statistic is $\sum_i a_i\, q_i(s_i)$, the total power $\|s\|^2$ for unit-variance Gaussian sources.
The rank of $L$ is a property of the segment covariances of the recording, which is what makes this condition testable without ground truth (Appendix~\ref{app:placement}).
